% Pursuing Stacks (cotes 134-2 to 134-8): archive cote 134-5, site page 38 (PDF page 39), typescript page 225.
% typescript. Draft compared: ps4.tex:1851-1911.
% AI-transcribed by Claude Opus 5.5 (claude-opus-5-5, adaptive thinking, effort high), Message Batches API, batch ps2 (msgbatch_01NVwr1Yuo8qJf5rykeGBXPK), 2026-11-12; not yet checked by eye.
\scanpage{134-5/38}
\mspage{225}

\[ \underline{Asph}(A,B)^{\circ}\hand{\longrightarrow}\underline{Hom}(Hot_{B},Hot_{A}) \quad , \qquad i \hand{\longmapsto} \overline{i^{*}} \quad , \]
where \uline{Asph} denotes the full subcategory of the functor category \uline{Hom}
made up with aspheric functors. This functor transforms arbitrary
arrows from the left hand side into isomorphisms on the right, and
therefore it factors through the fundamental groupoid (i.e. localization
of the left hand category with respect to the set of \uline{all} it's arrows):
\[ (\Pi_{1}(\underline{Asph}(A,B)))^{\circ} \hand{\longrightarrow} \underline{Hom}(Hot_{B},Hot_{A}) \quad . \]
If we now remember that we had assumed A and B to be test categories
(otherwise the functors just written would be still defined, but their
values would not necessarily be \uline{equivalences} between $Hot_{B}$ and $Hot_{A}$ ,
but merely functors between these), we may hope that this might imply
that the category $\underline{Asph}(A,B)$ of aspheric functors from A \inserted{whenever non-empty} to be \deleted{is}
1-connected. Whenever this is so, in any case, we get ``a priori''
(namely without any reference to Hot itself) a transitive system of
isomorphisms between the functors $\overline{i^{*}}$ , for i in $\underline{Asph}(A,B)$, \xout{Axxaxmaxxxr}
\xout{afxfaxx} hence a \inserted[typed]{``canonical''} \uline{functor} $Hot_{B}\hand{\longrightarrow}Hot_{A}$\hand{\unsure{)|,}} \uline{defined up to}\hand{\unsure{,}} \uline{unique isomorphism}
(and which, in case A and B are test categories, or more generally
pseudo-test categories\xout{*}, is an equivalence, and the one precisely
\margin{(1)} stemming from the diagram (\over{*}{1})).

\uline{Remark} Here the reflection slipped, almost against will, into a rela-%
ted one, about comparison of $Hot_{A}$ and $Hot_{B}$ for arbitrary A and B
(not necessarily test categories nor \uline{even} pseudo test categories), using
aspheric functors $i: A \hand{\longrightarrow} B$ to get $i^{*} : Hot_{B}\hand{\longrightarrow}Hot_{A}$ \over{,}{(} not necessa-%
rily an equivalence\hand{)}. As seen above, this functor depends rather loose-%
ly on the choice of i, and \xout{xxx} we could develop comprehensive conditions
on A and B, not \xout{xx} at all of a test-condition nature, implying that
\xout{xxxxx} just using the isomorphims $\overline{u}$ between these \inserted{functors}, we get a canonical
\uline{transitive} system of isomorphisms between the\unsure{m}, \deleted{namely} \inserted{hence} a \uline{canonical}
functor $Hot_{B}\hand{\longrightarrow}Hot_{A}$ , not depending on the \inserted{particular} choice of any aspheric
functor $i: A\hand{\longrightarrow} B$. As we are mainly interested in the modelizing case
though, I dont think I should dwell on this much longer here. Anyhow,
in case both A and B are pseudo-test categories, and provided only
$\underline{Asph}(A,B)$ is 0-connected (not necessarily 1-connected), it follows \xout{frm}
\margin{(1)} from comparison with the diagram (\over{*}{1}) above that the isomorphisms $\overline{u}$
do give rise \uline{ind}eed to a \uline{transitive} system of isomorphisms between the
equivalences $\overline{i^{\hand{*}}}$ .

In most cases though, such as $A =\hand{\Delta}$ and $B =\hand{\Box}$ say (the test

\margin{38}
